\section{The conditional geometric-epoch tradeoff}
\label{sec:epoch-blueprint}

The cleanest algebraic expression of volume flattening uses geometrically
decreasing tolerances
\begin{equation}
  \tau_j=2^{-j}\tau_0,
  \qquad J=\Theta(\log(\tau_0/\tau_J)).
  \label{eq:geometric-tolerances}
\end{equation}

\begin{proposition}[Conditional geometric-epoch work]
\label{prop:conditional-epochs}
Suppose that at epoch $j$ an algorithm has a certified region $U_j$ with
\begin{equation}
  \vol(U_j)\leq\frac{C}{\tau_j},
  \label{eq:epoch-volume}
\end{equation}
and that $H=C_0/\sqrt\alpha$ restricted iterations reduce the relevant error
by a constant factor.  If discovery and transition between epochs introduce
no additional full scans, then
\begin{align}
 T&=\bigO\!\left(\frac1{\sqrt\alpha}
                  \log\frac{\tau_0}{\tau_J}\right),
 \label{eq:epoch-iterations}\\
 \mathcal W
 &\leq H\sum_{j=0}^{J}\vol(U_j)
 =\bigO\!\left(\frac1{\sqrt\alpha\,\tau_J}\right).
 \label{eq:epoch-work}
\end{align}
\end{proposition}

\begin{proof}
The iteration bound is $H(J+1)$.  Because $1/\tau_j$ grows geometrically,
$\sum_j1/\tau_j=\bigO(1/\tau_J)$, which gives \cref{eq:epoch-work}.
\end{proof}

The proposition is only an accounting implication.  It assumes precisely the
two properties that need proof: a locally discoverable certified region and
stable acceleration as that region grows.
